\section{Compact fibres and the affine factor}
\label{cmp:section}

We retain the notation of Theorem~\ref{quo:projection} and
Theorem~\ref{grp:kernel-abelian}.  Thus $U=\widetilde X$, the base $M$ is
a bounded symmetric domain, and the image $\Lambda$ of $\Gamma$ acts
freely and properly discontinuously on $M$, with compact quotient.
The kernel $K$ is finitely generated and virtually abelian.  Passing to
the finite cover used to obtain this free action changes neither $U$
nor the assertion to be proved.  Here and below $\Gamma$ and $\Lambda$
denote the groups after this replacement, so $U/\Gamma$ is the chosen
finite cover of $X$.

We first construct, for some integer $a\geq0$, a proper map from $U$
to $M\times\mathbb C^a$ whose fibres are simply connected normal
projective varieties.  This step passes through relative Albanese
tori, which need not be
isotrivial.  We then prove that the compact fibres are all isomorphic
and that their family is a holomorphic product.

\subsection{The fibre cover and relative Albanese map}

We record precisely which regularity conclusions of the projection
construction are needed here.  Equivariance identifies the pullback
of $U/\Gamma\to M/\Lambda$ to $M$ with
\[
       U/K\simeq (U/\Gamma)\times_{M/\Lambda}M\longrightarrow M.
\]
This proper family and its connected finite covering families are flat and
topologically locally trivial.  Every fibre is connected, normal and
projective.  The simultaneous resolution gives a smooth proper
projective family over
$M$ mapping to each such covering family, with connected point inverses
and fibrewise birational maps.  These maps induce the pullback maps of
the integral first-cohomology local systems.  Finally, the universal
cover of every fibre has a semialgebraic projective-open presentation, up to
biholomorphism.  These are conclusions at \emph{every} point of $M$.
Normality of the total space, or generic normality alone, would not
suffice for the argument below.

\begin{lemma}[The characteristic fibre cover]
\label{cmp:characteristic}
There is a characteristic finite-index subgroup $K'\subset K$ which
is free abelian.  The quotient $\mathcal X=U/K'$ is normal and its map
$p:\mathcal X\to M$ is proper and flat.  It is topologically locally
trivial with connected normal projective fibres $N_m$, and
\[
          \pi_1(N_m)\xrightarrow{\ \sim\ }\pi_1(\mathcal X)=K'.
\]
The restriction of $U\to\mathcal X$ to $N_m$ is its connected universal
cover.
\end{lemma}

\begin{proof}
Choose a finite-index abelian subgroup of $K$.  It is finitely generated,
and hence contains a finite-index torsion-free subgroup $B$.  Put
$d=[K:B]$, and intersect all subgroups of $K$ of index at most $d$.
There are finitely many such subgroups, since homomorphisms from a
finitely generated group to each of the finite symmetric groups form
a finite set.  Their intersection $K'$ is characteristic, has finite
index, and is contained in $B$.  It is therefore free abelian.
Since $K$ is normal in $\Gamma$, so is $K'$.

The map $\mathcal X\to U/K$ is a finite covering.  Thus the stated
properness and regularity follow from the projection contract.  As
$U$ is simply connected, its quotient by $K'$ has fundamental group
$K'$.  A locally trivial bundle over the contractible manifold $M$
is a homotopy equivalence on each fibre; alternatively, its homotopy
exact sequence gives the displayed isomorphism.  The covering induced
on a fibre consequently corresponds to the zero subgroup of its
fundamental group, and is connected.
\end{proof}

\begin{lemma}[Relative Albanese descent]
\label{cmp:albanese}
There is an integer $a\geq0$ with $K'\simeq\mathbb Z^{2a}$ and a smooth proper family of
Albanese torsors $q:\mathcal A\to M$ and a proper morphism
\[
                 \alpha:\mathcal X\longrightarrow\mathcal A
\]
over $M$ whose restriction to $N_m$ is its Albanese morphism.  Both
constructions are equivariant for $\Gamma/K'$.  The morphisms
$\alpha|_{N_m}$ and $\alpha$ induce isomorphisms on fundamental groups.
\end{lemma}

\begin{proof}
Let $r:\mathcal Y\to\mathcal X$ be the simultaneous resolution and
let $p^+:\mathcal Y\to M$ be its smooth proper projective structure
map.  First, for each $m$, the map $r_m:Y_m\to N_m$ is surjective on
fundamental groups.  Indeed, lift it to the connected covering of
$N_m$ corresponding to the image of $\pi_1(Y_m)$.  Each connected
point inverse of $r_m$ maps to one point of the discrete covering
fibre.  Since a proper surjection is a quotient map, the lift descends
to a continuous section of this covering.  A connected covering with
a section has one sheet.  This proves the assertion.

It follows that
\begin{equation}
 W_{\mathbb Z}:=R^1p_*\mathbb Z\ \hookrightarrow
                         R^1p^+_*\mathbb Z=:H_{\mathbb Z}
 \label{cmp:cohomology-inclusion}
\end{equation}
is an inclusion of local systems, with saturated image.  Here the
local-system assertion uses the topological local triviality specified
above, and its compatibility with $r$ uses the map of the actual
families.  Saturation also follows directly: a homomorphism
$\pi_1(Y_m)\to\mathbb Z$ whose nonzero integral multiple vanishes on
$\ker(r_{m*})$ itself vanishes on that kernel.

We use the normal-projective Albanese theorem in its integral form:
if $N$ is normal and projective, then
\[
 H_1(N,\mathbb Z)\longrightarrow H_1(\operatorname{Alb}(N),\mathbb Z)
\]
is surjective with finite kernel
\cite[Lemma~4.7]{ClaudonHoeringKollar2013}.  In the present situation
$H_1(N_m,\mathbb Z)=K'$ is torsion-free, so this map is an isomorphism.
In particular its rank is $2a$.  Moreover, the image of
$H^1(N_m,\mathbb Q)$ in $H^1(Y_m,\mathbb Q)$ is a Hodge substructure:
it is exactly the pullback from $\operatorname{Alb}(N_m)$ under the
algebraic map $Y_m\to N_m\to\operatorname{Alb}(N_m)$.  This observation
also derives the required purity without an assumption about the
singularities being rational.

Let $F^1H_{\mathcal O}$ be the holomorphic Hodge subbundle of the
weight-one variation of the smooth projective family $p^+$.
Holomorphicity holds over this analytic base by Griffiths's
period-mapping theorem \cite[II, Theorem~(1.1)]{Griffiths1968}.  The
intersection
\[
                    W^{1,0}=W_{\mathcal O}\cap F^1H_{\mathcal O}
\]
is a holomorphic subbundle: it is the kernel of the holomorphic map
$W_{\mathcal O}\to H_{\mathcal O}/F^1H_{\mathcal O}$ and has the
constant rank $a$.  Together with $W_{\mathbb Z}$ it defines a
polarizable weight-one variation.  The lattice dual to $W_{\mathbb Z}$
embeds by evaluation in
$E=(W^{1,0})^\vee$, and the quotient gives a smooth family $E/W_{\mathbb Z}^\vee$
of abelian varieties, by the integral analytic Jacobian construction
\cite[\S2.1]{Claudon2018}. Fibrewise the lattice is a full lattice;
continuity in a local basis makes the quotient a holomorphic torus
family, with no change of lattice index.

On a small base ball, choose a section of the smooth family $p^+$.
The relative Albanese map of $\mathcal Y$ based at that section,
followed by the quotient specified by $W$, is a holomorphic map to
this torus family; see the smooth construction in
\cite[\S1.4]{Fujiki1983} and
\cite[II, \S2(a) and Example~(2.15)]{Griffiths1968}.  Equivalently, integrate the relative holomorphic
one-forms in $W^{1,0}$.  Its map on integral first homology is the
quotient induced by $r_m$, so its restriction to every point inverse
of $r_m$ has trivial induced fundamental-group map.  Each compact
connected point inverse therefore lifts to the vector-space cover of
the target torus.  Holomorphic functions on a compact connected
reduced complex space are constant, so this restriction is constant.

Thus the relative map is constant on the fibres of $r$.  Since
$\mathcal X$ is normal and $r$ is a proper modification,
$r_*\mathcal O_{\mathcal Y}=\mathcal O_{\mathcal X}$; the map descends
holomorphically to $\mathcal X$.  Changing the chosen local section
changes only the origin in the torus family.  The resulting local
maps glue to the family of Albanese torsors $\mathcal A$ and to
$\alpha$, as in the translation-cocycle description of
\cite[\S2.2, Proposition~2.1]{Claudon2018}.  The fibrewise universal property identifies this map with
the normal Albanese morphism, which is algebraic on every projective
fibre.  An element of $\Gamma/K'$ acts on the intrinsic system
$W_{\mathbb Z}$ and its Hodge subbundle, and thus induces a
holomorphic linear map of the torus families.  In charts with origins,
the additional translation is obtained by evaluating the relative
Albanese map at the image of the chosen local section; it is
holomorphic.  The fibrewise universal property makes these maps agree
on overlaps and obey the group law.  This gives the asserted
$\Gamma/K'$-equivariance of the torsor construction.

Properness of $\alpha$ follows from properness of $p$: a compact
subset of $\mathcal A$ projects into a compact subset of $M$, and its
inverse image is closed in the compact inverse image under $p$.
The fibrewise fundamental-group assertion follows from the integral
Albanese statement and $\pi_1(N_m)=K'$.  The smooth proper torus family
$q$ is topologically locally trivial.  Since $M$ is contractible, its
fibres and those of $p$ induce isomorphisms on fundamental groups.
The assertion for $\alpha$ follows from the commutative diagram of
these two fibre inclusions.
\end{proof}

\subsection{The universal affine base}

\begin{proposition}[The proper map to the affine base]
\label{cmp:affine-base}
There is a $\Gamma$-equivariant proper flat surjection
\[
                    j:U\longrightarrow T,
          \qquad T\simeq M\times\mathbb C^a,
\]
whose fibres are connected, simply connected, normal projective
varieties.  The action of $\Gamma$ on $T$ is proper and cocompact and
has finite stabilizers.  A polarization of the finite cover of $X$
pulls back to a $\Gamma$-invariant line bundle $L$ on $U$, ample on
each fibre of $j$.
\end{proposition}

\begin{proof}
We apply the following precise form of the Koll\'ar--Pardon theorem
\cite[arXiv version~2, Theorem~20]{KollarPardon2012}: if $N$ is normal projective, $B$
is smooth projective with contractible universal cover, and
$g:N\to B$ is a morphism such that the \emph{entire} pullback to that
cover is biholomorphic to a semialgebraic open subset of a projective
variety, then $g$ is a holomorphic fibre bundle.  Surjectivity is part
of its conclusion; connected fibres are not a hypothesis or an
automatic part of the general statement.

Apply this to $N_m\to\operatorname{Alb}(N_m)$.  The target is an
abelian variety with universal cover $\mathbb C^a$.  By
Lemma~\ref{cmp:albanese}, the pullback cover is connected and is
precisely the universal cover of $N_m$.  The required semialgebraic
presentation is therefore the one supplied by the projection
construction.  Each Albanese map is a holomorphic fibre bundle and
is flat.  Its fibres are connected: in the homotopy sequence of a
bundle, surjectivity on $\pi_1$ implies that $\pi_0$ of the fibre is
a singleton.  Since $\pi_2$ of a torus is zero and the map on $\pi_1$
is an isomorphism, the same sequence shows that each fibre is simply
connected.  Local holomorphic products with a ball and normality of
$N_m$ show that the fibre is reduced and normal.  It is projective,
being an algebraic fibre of a morphism of projective varieties.  These
arguments include $a=0$, when the Albanese map is the map to a point.

The fibrewise flatness criterion now applies to
$\mathcal X\to\mathcal A\to M$: the source is flat over $M$, the
target is smooth over $M$, and the maps on all fibres over $M$ are
flat.  More explicitly, for the analytic local rings $R\to B\to C$
at $m$, $\alpha(x)$ and $x$, respectively, $B$ and $C$ are flat over
$R$, and $C/\mathfrak m_R C$ is flat over
$B/\mathfrak m_R B$.  The fibrewise flatness criterion for Noetherian local rings
\cite[Tag~00MP]{StacksProject} applies with module $C$, which is
finite free over itself, and implies that $C$ is flat over $B$.  Thus $\alpha$ is flat
at every point.

Let $T\to\mathcal A$ be its universal covering.  The pullback
$\mathcal X\times_{\mathcal A}T$ is connected and simply connected,
by the isomorphism on fundamental groups in Lemma~\ref{cmp:albanese}.
It is consequently $U$, and its projection $j$ is proper and flat by
base change.  It is surjective by the fibrewise surjectivity just
proved.  Its fibres have all the asserted properties.

We explain the structure of $T$ without requiring the torus family
to be isotrivial.  On a small contractible ball $B\subset M$, choose
an origin for $\mathcal A_B$.  Relative exponential identifies its
fibrewise universal cover with the vector bundle $E_B$.  The inclusion
of a fibre in $\mathcal A$ induces an isomorphism on $\pi_1$, since
$M$ is contractible.  Hence the restriction of the global universal
cover $T$ to $\mathcal A_B$ is connected and is exactly this cover
$E_B\to\mathcal A_B$.  On overlapping balls, changes of origin lift
to translations, while the linear terms are the transition maps of
the relative Lie bundle $E$.  Thus $T\to M$ is an affine bundle with
linear part $E$.  Its translation cocycle lies in
$H^1(M,\mathcal O(E))=0$ by Cartan's theorem~B
\cite[Expos\'e~18, \S4]{Cartan1952}.  It follows that
$T\simeq E$.  The vector bundle $E$ is topologically trivial because
$M$ is contractible, and holomorphically trivial by the Oka principle
\cite{Grauert1958,ForstnericPrezelj2000}.
Thus $T\simeq M\times\mathbb C^a$, which is Stein, contractible and
homeomorphic to Euclidean space.

\enlargethispage{\baselineskip}
The functorial action on $\mathcal A$ lifts to $T$, since $T$ is its
universal cover.  Choose a lift $\gamma_T$ whose value at one point
agrees with $j\circ\gamma$.  Then $\gamma_T\circ j$ and
$j\circ\gamma$ are lifts from the connected space $U$ of the same
map to $\mathcal A$ and agree at one point, so they agree everywhere.
Surjectivity of $j$ then gives uniqueness of $\gamma_T$ and the
group law.  If a compact
$C\subset T$ meets $\gamma C$, then the compact set $j^{-1}(C)$ meets
its $\gamma$-translate.  The deck action on $U$ is proper, so only
finitely many $\gamma$ have this property.  This proves properness
on $T$, including finite stabilizers and a possible finite kernel.
The induced surjection $U/\Gamma\to T/\Gamma$ shows that the latter
quotient is compact.  Finally, pull back an ample line bundle from
$U/\Gamma$.  It is $\Gamma$-invariant.  On a compact fibre its map to
$U/\Gamma$ is finite onto its image, being proper and locally
injective, so this pullback is ample on the fibre.
\end{proof}

\subsection{Variation of the compact fibres}

It remains to show that the compact fibres of $j$ are all
polarized-isomorphic.  We shall put their parameter space $T$ in a
projective Chow space and show that each polarized isomorphism class
is a semialgebraic, $\Gamma$-invariant subset.  The following
topological lemma will force its closure to fill $T$; the
semialgebraic description will then rule out two distinct dense
classes.  Once all fibres are polarized-isomorphic, relative Hilbert
embeddings and the Oka principle give a holomorphic product.

\begin{lemma}[Rigidity for invariant closed subsets]
\label{cmp:closed-rigidity}
Let a countable discrete group act properly and cocompactly by
diffeomorphisms on a contractible smooth manifold $T$ homeomorphic to
$\mathbb R^b$.  If a nonempty closed invariant subset $A\subset T$
has the homotopy type of a finite CW complex, then $A=T$.
\end{lemma}

\begin{proof}
We give the controlled proper-homotopy argument, including the case
of finite stabilizers; compare \cite[\S1, Lemmas~11--12 and
Proposition~13]{KollarPardon2012}.  A proper cocompact smooth action
admits an invariant complete Riemannian metric.  Its distance $d$ is
proper.  Fix $o\in T$ and $R>0$ such that the balls $B(\gamma o,R)$
cover $T$.  Choose continuous functions $\lambda_\gamma:T\to[0,1]$
equal to one on these balls and zero outside $B(\gamma o,R+1)$,
obtained by translating one function of distance.  Properness makes
their supports locally finite.  Cocompactness and properness also
give a uniform bound $N$ on the number of nonzero terms at any point.
Indeed translate that point into a fixed compact set and count orbit
centres in its compact $(R+1)$-neighbourhood.  Finite stabilizers
merely give finite repetitions in this count.

Choose a homotopy $H:A\times[0,1]\to A$ with $H_0=\operatorname{id}$
and $H_1(A)$ contained in a compact subset $C\subset A$.  Such a
homotopy is obtained by factoring a homotopy equivalence through a
finite CW complex.  Order the countable group once and for all.
For $y\in T$, apply, in that order, the finitely many maps
\[
       a\longmapsto\gamma H(\gamma^{-1}a,\lambda_\gamma(y))
\]
to a fixed point $a_0\in A$.  Call the resulting point $s(y)$.
Local finiteness and $H_0=\operatorname{id}$ imply that $s$ is
continuous, including where a factor enters or leaves its support.

At least one factor has time one.  Immediately after that factor the
point lies in $\gamma C$, at uniformly bounded distance from $y$.
Subsequent factors preserve a uniform distance bound.  In detail,
if $d(a,y)\le B$ and $\lambda_\delta(y)\ne0$, then
$d(\delta^{-1}a,o)\le B+R+1$.  The homotopy of the compact set
$A\cap\overline B(o,B+R+1)$ has compact image, so the distance of
$\delta H(\delta^{-1}a,t)$ from $y$ has a bound depending only on
$B$.  Iterating at most $N$ times gives
\[
                        d(s(y),y)\le B_0
\]
for a fixed $B_0$.  In particular $s:T\to A$ is proper.

For completeness, bounded distance from the identity gives a proper
homotopy in this situation as follows.  Choose a contraction
$J:T\times[0,1]\to T$ from the identity to $o$.  Form $Q(y,x,t)$
by composing the translated contractions
\[
       x\longmapsto\gamma J(\gamma^{-1}x,t\lambda_\gamma(y))
\]
in the same fixed order.  At $t=1$ a full contraction occurs, so the
output is independent of $x$.  If $d(x,y)$ is bounded, all the
intermediate outputs stay at a uniformly bounded distance from $y$;
the compact-set argument above applies now also before the first
full contraction.  Join $y$ to $Q(y,y,1)$ by this construction and
return along the construction starting at $s(y)$.  The two middle
points agree.  This is a continuous proper homotopy from
$\operatorname{id}_T$ to the composite $T\xrightarrow{s}A\hookrightarrow T$.

If $b=0$ the conclusion is immediate.  Otherwise, if $x\notin A$,
the latter proper map omits $x$.  Its extension to the one-point
compactification $T^+\simeq S^b$ therefore has degree zero, whereas
the proper homotopy to the identity gives degree one.  This
contradiction proves the lemma.  No assertion that $T\to T/\Gamma$
is a covering has been used.
\end{proof}

\begin{lemma}[The semialgebraic space of fibres]
\label{cmp:chow-base}
Fix a semialgebraic realization $U\subset Z$, where $Z$ is projective,
and a projective embedding of $Z$.  Let $n$ be the fibre dimension of
$j$ and $d$ the degree of one fibre.  In the projective Chow space
$\operatorname{Chow}_{n,d}(Z)$, the locus
\[
 C=\{c:\ c\text{ is an integral cycle and }|c|\subset U\}
\]
is semialgebraic, and the map $t\mapsto[j^{-1}(t)]$ is a homeomorphism
from $T$ onto $C$.
\end{lemma}

\begin{proof}
A proper flat family of pure-dimensional compact subspaces over a
reduced complex base defines an analytic family of fundamental cycles;
see \cite[p.~27, Remark~(iv)]{Barlet1999}. Its map to the projective
Chow space is continuous. Indeed, view cycles in $Z$ as cycles in
a fixed smooth projective space. Support and integration continuity
\cite[\S2.4(b), Lemma~4 and \S2.5, Theorem~4]{AndreottiNorguet1967}
identify the usual Chow topology with the analytic cycle topology:
the natural map is a continuous bijection from the compact
fixed-degree Chow variety onto the corresponding fixed-degree locus
in the Hausdorff cycle space, hence a homeomorphism. This comparison
uses no smoothness of $Z$.  Apply this to the graph of
$j$ in $T\times Z$.  Every fibre is connected and normal, hence
irreducible and reduced, so its fundamental cycle is integral.  The
degree is locally constant in this family, and hence equals $d$ on
the connected base $T$.

The integral-cycle locus is constructible.  More explicitly, its
complement is the finite union, over $d_1+d_2=d$ with both degrees
positive, of the images of the proper addition maps from the
corresponding Chow spaces.  These images include nonreduced cycles.
Containment of the support in $U$ is a semialgebraic condition by the
algebraic universal incidence and real quantifier elimination.
Thus $C$ is semialgebraic.

Any compact irreducible analytic subvariety of $U$ maps to a point of
the Stein manifold $T$: holomorphic functions are constant on that
subvariety, and holomorphic functions on a Stein manifold separate
points.  An integral $n$-cycle in $U$ therefore lies in a fibre and,
by equality of dimension and irreducibility of that fibre, is its
whole reduced cycle.  This proves that the displayed cycle map is
bijective.  Its continuity follows from the analytic-family statement.
For inverse continuity, if fibre cycles $c_\nu$ converge to the cycle
of $j^{-1}(t)$, choose a regular point $x$ of that fibre.  Convergence
of cycles supplies points $x_\nu\in|c_\nu|$ tending to $x$; for
example, use a small transverse disk at $x$, whose local intersection
number is one.  Then their parameters satisfy
$j(x_\nu)\to j(x)=t$.  All spaces involved are metrizable, so this
proves inverse continuity.
\end{proof}

\begin{lemma}[Integral marking of polarized isomorphism classes]
\label{cmp:polarized-locus}
For every fixed polarized fibre $(F,L_F)$ of $j$, the locus of
fibres polarized-isomorphic to it is semialgebraic in $C$.
\end{lemma}

\begin{proof}
Write $A_Z=\mathcal O_Z(1)|_U$ and choose $k>0$ for which
$H=L_F^k$ is very ample.  If $g:F\to j^{-1}(t)$ is a polarized
isomorphism, its graph in $F\times Z$, with the product projective
embedding defined by $H\boxtimes\mathcal O_Z(1)$, has degree
\begin{equation}
 D=\int_{j^{-1}(t)}
                \bigl(kc_1(L)+c_1(A_Z)\bigr)^n.
 \label{cmp:graph-degree}
\end{equation}
This integer is independent of $t$.  Here is a direct way to check
the constancy, which does not extend $L$ to $Z$.  Choose a smooth
Hermitian metric on $L$ on the complex space $U$, and use its
curvature form, together with a Fubini--Study form for $A_Z$.  Near
any fixed parameter, properness confines the supports of the fibre
cycles to a compact subset of $U$.  Integration of the resulting
smooth $2n$-form over these cycles is continuous in their analytic
cycle family.  Each value is the integral intersection number in
the displayed formula, so it is locally constant.  Smooth forms on
a complex space here mean forms locally extended from an ambient
manifold; a finite partition of unity on the compact set gives the
usual continuity-of-integration test.  The Chern integral equals the
Cartier intersection number also for a singular fibre: pull the
line bundles back to a projective resolution and use its degree-one
pushforward of the fundamental cycle.  The same argument proves
constancy of each mixed intersection separately.

Consider the integral-cycle locus in the projective Chow space of
$n$-cycles of degree $D$ in $F\times Z$, and let $P$ be its locus of
cycles $G$ such that
\[
 |G|\subset F\times U,\qquad
 p_F:|G|\to F\text{ is bijective},\qquad
 p_Z:|G|\to U\text{ is injective}.
\]
All three conditions are first-order conditions on the universal
incidence, and hence are semialgebraic.  The first projection is a
proper quasi-finite algebraic map, hence finite.  It is birational,
since in characteristic zero a finite map bijective on complex
points has degree one.  Its target $F$ is normal, so it is an
isomorphism.  The second projection has an irreducible compact image
of dimension $n$, which by the preceding lemma is a whole fibre of
$j$.  This image is normal, and the same finite birational argument
shows that the second projection is an isomorphism onto that fibre.
Consequently $P$ parametrizes exactly the isomorphisms from $F$ to
fibres whose graphs have the specified degree.

The universal incidence over $P$ maps bijectively and closedly to
$F\times P$: closedness follows by projection along the compact
factor $Z$.  It is therefore homeomorphic to $F\times P$.  In
particular it gives a continuous evaluation map
\[
                         e:F\times P\longrightarrow U.
\]
Connected semialgebraic sets are path connected, and $P$ has only
finitely many connected components.  A path in one such component
gives a homotopy of the maps $e_p:F\to U$.  Hence
$c_1(e_p^*L)\in H^2(F,\mathbb Z)$ is constant on each component.
This is the integral marking test; neither definability of $L$ nor
an extension of $L$ to the compactification has been assumed.

For a simply connected normal projective variety $F$, the map
\[
             c_1:\operatorname{Pic}(F)\longrightarrow H^2(F,\mathbb Z)
\]
is injective.  Indeed $\operatorname{Pic}^0(F)$ is projective for
normal projective $F$ and is smooth in characteristic zero
\cite[Theorem~5.4 and Remark~5.6]{Kleiman2005}.  GAGA identifies algebraic and analytic line bundles on $F$
\cite[Proposition~18]{Serre1956}. The exponential
sequence and $H^1(F,\mathbb Z)=0$ identify $\operatorname{Pic}^0(F)$ analytically
with the vector group $H^1(F,\mathcal O_F)$.  A compact complex vector
group is zero.  The kernel of $c_1$, which is the image of this
vector group in the exponential sequence, is therefore zero.

It follows that the condition $e_p^*L\simeq L_F$ selects a union
$P_L$ of connected components of $P$, so $P_L$ is semialgebraic.
Taking the second image of a graph gives a semialgebraic map to $C$.
One can see this directly by expressing equality of supports as
\[
 z\in|c|\quad\Longleftrightarrow\quad
              (\exists x\in F)\ (x,z)\in|G|;
\]
integrality makes this support determine the cycle.  The image of
$P_L$ is exactly the requested polarized isomorphism class, including
all its members by the graph-degree calculation.  Quantifier
elimination proves the assertion.
\end{proof}

\begin{proposition}[Triviality of the compact factor]
\label{cmp:product}
There is a simply connected normal projective variety $F$ such that
\[
                    \widetilde X\simeq M\times\mathbb C^a\times F.
\]
\end{proposition}

\begin{proof}
By Lemma~\ref{cmp:polarized-locus}, every polarized isomorphism class
is a semialgebraic subset of the semialgebraic model $C$ of $T$.
It is $\Gamma$-invariant, since the polarization is pulled back from
$U/\Gamma$.  Its closure in $C$ is a nonempty closed invariant
semialgebraic subset.  Compatible semialgebraic triangulation
\cite[\S3, Theorem~4]{Lojasiewicz1964} gives that closure the
homotopy type of a finite CW complex.  Under the homeomorphism
of Lemma~\ref{cmp:chow-base}, Lemma~\ref{cmp:closed-rigidity} therefore
forces it to be all of $T$.  Thus every class is dense.  Two disjoint
semialgebraic subsets cannot both be dense in the same semialgebraic
space: finite cell decomposition shows that a dense semialgebraic
subset contains a dense relatively open subset.  All fibres are
consequently isomorphic as polarized varieties.

We next justify local holomorphic triviality for these possibly
singular fibres.  Choose $r$ so large that $L_F^r$ is very ample and
$H^i(F,L_F^r)=0$ for $i>0$.  The same holds on every fibre.  The constant-cohomology base-change theorem
\cite[\S7, Satz~5]{Grauert1960} makes $j_*L^r$ locally free, with
its formation commuting with fibres: $L^r$ is flat over the base
and all its fibre cohomology dimensions are constant.  After trivializing it on a
small open set $B\subset T$, the evaluation map embeds
$j^{-1}(B)$ in $B\times\mathbb P^q$ and gives a holomorphic Hilbert
map $h:B\to\operatorname{Hilb}(\mathbb P^q)$.
Here one uses the analytic-base formulation: Douady's universal
proper-flat family \cite[\S\S9.7--9.8, Theorems~1--2]{Douady1966} agrees
with the analytification of the projective Hilbert scheme
\cite[\S3, Theorem~3.2]{Grothendieck1961}. The natural map
between them has the same points and the same local deformation
functors, since GAGA applies on every Artinian complex base
\cite[Expos\'e~XII, Theorem~4.4]{SGA1};
thus it is an isomorphism of analytic germs. This supplies the
holomorphic classifying map also for singular fibres.  All its points lie
in the single orbit of the embedded $F$ under
$G_0=\operatorname{PGL}_{q+1}(\mathbb C)$.

This orbit is a locally closed smooth complex subvariety.  Since
$B$ is reduced, a holomorphic map with image in it factors
holomorphically through it.  For its closed algebraic stabilizer
$G\subset G_0$, the quotient map $G_0\to G_0/G$ has local
holomorphic sections, by the holomorphic submersion theorem.
Locally lifting $h$ and applying the inverse projective
transformations gives a holomorphic product $B\times F$.
The resulting transition maps take values in the complex linear
algebraic group $G$.  Thus $j$ is the bundle associated to a
holomorphic principal $G$-bundle on $T$.

Since $T$ is contractible and paracompact, that principal bundle is
topologically trivial.  The Oka principle for complex Lie groups
over a Stein base makes it holomorphically trivial; one may use
\cite[Theorem~1.3 and Example~(A)]{ForstnericPrezelj2000} on its
continuous trivializing section.  If necessary, first reduce to
the identity component of $G$ using the triviality of the discrete
covering $G/G^\circ$ over the simply connected space $T$.
We obtain $U\simeq T\times F$.  Proposition~\ref{cmp:affine-base}
gives the required product, with all the stated properties of $F$.
\end{proof}

\begin{proposition}[The converse]
\label{cmp:converse}
If $\widetilde X\simeq D\times\mathbb C^m\times F$, with $D$ a
bounded symmetric domain and $F$ projective, then $\widetilde X$ is
biholomorphic to a semialgebraic open subset of a projective variety.
\end{proposition}

\begin{proof}
In the Borel embedding, $D$ is an open orbit in its projective
compact dual $D^\vee$ \cite[\S4, Theorem~2 and \S6, Proposition~2]{Borel1954};
see also \cite[\S2.1.1, Borel embedding theorem]{FalbelGuillouxWill2024}.
The acting adjoint real semisimple group is
the identity component of the real points of a real algebraic group
\cite[Proposition~1.7]{Milne2005}; it is semialgebraic, and its action on $D^\vee$ is algebraic.
The image of its orbit map is semialgebraic by real quantifier
elimination.  Thus $D\subset D^\vee$ is a semialgebraic open subset.
Taking the usual affine chart $\mathbb C^m\subset\mathbb P^m$ gives
\[
 D\times\mathbb C^m\times F
        \ \subset\ D^\vee\times\mathbb P^m\times F
\]
as a semialgebraic open subset of a projective variety.  Point
factors are allowed.  This construction concerns the biholomorphism
type of the covering space and imposes no algebraicity condition on
deck transformations.
\end{proof}
